%% =====================================================================
%%  T-01-04  Intent Is the Only Ethical Switch
%%  -------------------------------------------------------------------
%%  One idea per file. Core is mandatory. Slots in canonical order.
%%  Max 700 words. No layout commands. No filler. New source material
%%  goes in sources/<BOOK>/T-01-04.tex -- never by rewriting this file.
%% =====================================================================
%% @kind: topic
%% @id: T-01-04
%% @title: Intent Is the Only Ethical Switch
%% @subtitle: The same move is a defence or an attack depending on purpose
%% @chapter: c01
%% @order: 40
%% @status: stable
%% @sources: B1, B2
%% @updated: 2026-09-19
%% @allow-rewrite: slot migration to the seven-question format (docs/RULES.md section 2)

\topic{T-01-04}{Intent Is the Only Ethical Switch}{The same move is a defence or an attack depending on purpose}


\begin{core}
Nothing in the technique set distinguishes a defender from an operator. The same read, the same delay, the same controlled disclosure serves both. What separates them is the purpose the move is put to, and the only place that purpose is visible is in what the target is left holding afterwards.
\end{core}

\begin{mechanism}
A technique is a way of arranging a situation. Arrangement is neutral in the same sense that a lock is neutral: it keeps people out or keeps them in depending on which side you stand on. Learning the arrangement does not commit you to a side, but it does commit you to knowing which side you are on, because the moves feel identical from the inside. The practical test is external --- whether the other person ends up better informed or worse, whether they could reconstruct what happened if they tried, and whether you would be comfortable with them knowing exactly what you did and why.
\end{mechanism}

\begin{application}
  \item Decide the outcome you want before choosing a method, and check whether the method serves the outcome or merely the moment.
  \item Prefer methods that leave the other party's picture of the situation more accurate, not less.
  \item Where you must arrange a situation, arrange it toward a result you would defend aloud.
\end{application}

\begin{feedback}
  \item You are editing what someone knows in order to get a decision you have already made.
  \item You would not want the method described to the person you are using it on.
  \item The other party ends the exchange less able to see their own position than they were at the start.
\end{feedback}

\begin{failure}
Intent is not observable, including to yourself in the moment. People running a self-serving move generally experience it as reasonable, and the rationalisation arrives before the act rather than after. This entry is therefore a discipline rather than a guarantee --- useful because it forces the question before the move, weak because the answer it produces is often flattering. The external tests above are stronger than introspection for exactly that reason.
\end{failure}

\begin{countermeasures}
  \item Use the disclosure test: if the full method were described to the other person in advance, would the exchange still work? If not, it is manipulation rather than negotiation.
  \item Use the residue test: afterwards, is the other person better or worse able to look after their own interests?
  \item If you are the target, the equivalent question is whether you are being left holding a decision you cannot reconstruct.
\end{countermeasures}

\sources{B1, B2}
\seesources
\seealso{T-01-02, T-01-03, T-27-01}
